% ===========================================================================
% 04a-ruler.tex -- sec:res-metrics and sec:res-blind (legacy anchors kept).
% ===========================================================================

\beatsection{E}{1}{The metric is the first result}{sec:res-metrics}
\label{sec:methods-metrics}   % legacy anchor, so v1 cross-references resolve

\claimlede{The field's standard score is saturated by a predictor that carries
no information, and of the five metrics audited only \NIR{} grades in the
direction its name implies.}

\begin{question}[1]
Every claim in this chapter is a number handed down by a scoring rule, so the
scoring rule is what has to be audited first. The instrument grades metrics, not
models. It asks each candidate to separate a positive control from a predictor
that knows no drug, and four of the five candidates fail.
\end{question}

\noindent
A reported score is two claims at once, one about the model and one about the
rule that graded it. Nobody checks the second. Two questions decide it. Does the
score reward the thing it is named after, or can it be earned another way? And
does it register a true positive when there is one?

\paragraph{The score the field reports, and two companions} That score is a correlation between
predicted and true change from control, restricted to the genes that actually
move. The benchmarks of \cref{sec:lit-dispute} report it,
Miller et al.~\cite{miller} find it poorly calibrated, DrEval~\cite{dreval}
counts it among six standing obstacles. Write $x$ for the control profile, $y$
for the treated truth and $\hat y$ for the
prediction,\seealso{\cref{sec:methods-data} fixes $G$.} and let
\begin{equation}
 S_K \;=\; \operatorname*{arg\,top-}K_{\,g} \; \lvert y_g - x_g \rvert
\end{equation}
be the $K$ genes with the largest true change. Then
\begin{equation}
 \DEdr \;=\; \operatorname{corr}\!\big( \left(\hat{y} - x\right)\big|_{S_K}, \;
 \left(y - x\right)\big|_{S_K} \big).
\end{equation}
All three are the panel profiles converted to ranks, absent genes tied at the
worst rank $G$ (\cref{sec:app-convention}), so $S_K$ is selected on rank change
too; the \emph{expression frame} label used elsewhere names the object scored, a
cell profile rather than a residual, and not the units. We report $K = 50$,
Pearson; $K \in \{20, 100, 200\}$ and Spearman move the
numbers by under $0.02$, except at $K = 200$, where the top-$K$ set starts to
include genes that do not move.

Two companions to it are reported below, and neither is that score. Writing
$s = (y-x)|_{S_K}$ and $\hat{s} = (\hat{y}-x)|_{S_K}$ for the true and the
predicted shift, and $r$ for Pearson correlation on $S_K$,
\emph{partial-}$\DEdr$ regresses the control out of both arguments first:
\begin{equation}
 \textrm{partial-}\DEdr \;=\;
 \operatorname{corr}\!\big( \hat{s}, \, s \;\big|\; x|_{S_K} \big)
 \;=\;
 \frac{r_{\hat{s}s} \;-\; r_{\hat{s}x}\, r_{sx}}
      {\sqrt{\left(1 - r_{\hat{s}x}^{2}\right)\left(1 - r_{sx}^{2}\right)}}.
\end{equation}
Both arguments of $\DEdr$ carry the same $-x$, so a predictor that merely
reverts toward the control inherits agreement it did not earn; conditioning on
$x$ leaves only what a prediction has beyond control-reversion, and the quotient
is undefined rather than small when either shift is collinear with the control
on $S_K$, its denominator vanishing. The second companion moves the other lever.
$\paneltau$ is Kendall's $\tau$ between the predicted and the true rank vector
over all $G$ panel genes,
\begin{equation}
 \paneltau \;=\; \tau\!\left(\hat{y}, \, y\right),
\end{equation}
with no top-$K$ restriction and no control subtracted: where $\DEdr$ picks its
genes from the truth and scores a change, $\paneltau$ scores the ordering itself
on a gene set fixed before any prediction is read. It depends on the absent-gene
convention rather than escaping it, every omitted gene being tied at $G$, so it
registers omissions that $\DEdr$ sees only when an omitted gene falls in $S_K$,
and it is undefined when a prediction ties every gene.

\paragraph{What the metric actually rewards} Score a predictor that contains
nothing. Place every gene at the middle rank.\gloss{revert\_center}{Predicts
every panel gene at the panel's middle rank, defined from the metric alone and
needing no model, training set or atlas.} It scores $\DEdr(K{=}50) =$
\ci{0.999913}{0.999908}{0.999919}, above the two-real-cell positive control of
$0.76$ and above the model's $0.73$. Less information yields a higher score,
the hallmark of an artefact; \cref{fig:calibration}, panel a, sets the five arms
against that reference on one axis.

Its partial-$\DEdr$ and $\paneltau$ are undefined, not merely poor, for two
separate reasons. The prediction is the constant midpoint rank $\hat y = G/2$,
so the shift from control is $G/2 - x$, an exact affine function of $x$ from
which nothing survives partialling. And every predicted gene being tied, no rank
ordering exists whose Kendall $\tau$ can be computed. The zero-shift
construction, easily confused with this one, is \texttt{control\_copy},
$\hat y = x$; \texttt{revert\_center} is not it.

\begin{table}[htbp]
% MARGINALIA: in-column; caption and the one-line note in the margin.
\begin{lrbox}{\tvfb}
\begin{tabular}{@{}l l r l@{}}
\thead{Predictor} & \thead{Information} & \thead{$\DEdr$ (K50)} & \thead{partial-$\DEdr$} \\
\midrule
\arm{revert_center} (genes $\rightarrow G/2$) & none & \key{0.9999} & \quiet{undefined} \\
\arm{revert_mean} (predict mean control) & none, no fit & 0.961 & 0.25 \\
ridge linear (control $\rightarrow$ shift) & control & 0.947 & 0.28 \\
\addlinespace[10pt]
positive control (two real cells) & \quiet{--} & 0.76 & \quiet{--} \\
model (1B LLM) & control $+$ drug & 0.73 & \quiet{not measured} \\
\end{tabular}
\end{lrbox}
\sidecap{table}{\captiontitle{$\DEdr$ is saturated by a control artefact.} The two arms
needing no drug information sit above the positive control and above the model;
regressing the control out leaves $\approx 0.26$ of drug-agnostic skill in the
two rows where the partial is defined.\label{tab:de-artifact}
  \par\addvspace{6pt}The $\DEdr$ column is tier 2; the partial column, tier 1.}
\end{table}

\paragraph{Two innocent properties compose into the failure} $S_K$ is selected
using the truth, and treated
cells in this atlas regress toward a shared response, so a prediction reverting
toward the population centre correlates with $y - x$ by construction.
Partial-$\DEdr$ regresses the control out of both arguments first, closing that
route and collapsing the raw $0.95$ to $\approx 0.26$ in the two baselines with
a defined partial value, one of which performs no fitting at all.

The model's own partial-$\DEdr$ was never computed,\footnote{It requires a GPU
pass that was not run, so the cell in \cref{tab:de-artifact} is marked, not
filled with an expected value.} so the comparison anchors to $\paneltau$,
Kendall's $\tau$ over the full $\num{946}$-gene panel, measured for every arm.
The calibration  finds $\paneltau$ inverted, so neither an absolute claim
nor an ordering read off it counts as evidence of skill. What it can still show
is an \emph{absence} of separation. \texttt{revert\_mean} scores $0.268$, ridge
$0.265$ and the model $0.257$ (tier 2, \texttt{worst} convention), three arms
inside a spread of $0.011$.

The shift therefore decomposes into three parts: control reversion, an
artefact; a generic drug programme worth $\approx 0.26$, trivially matched; and
a drug-specific component no arm is shown capturing. Later residual-frame
instruments show the checkpoint reads trained-drug information while still
failing condition-specific prediction.

\paragraph{The spike-in, in brief} One instrument recurs from here on and is
worth naming before it is used. The \emph{spike-in} is a forced choice between
two drug populations measured within a plate, and no model enters it: a
reference pseudobulk is drawn from drug A's cells, and a metric must score a
disjoint sample of A's own cells closer to that reference than a sample of drug
B's, chance $0.50$. Difficulty is then dialled by titration --- B's population is
contaminated with a known fraction of A's cells --- so the task runs from easy at
zero contamination to impossible once B has been replaced by A and the two
candidates are one population. Because the right answer is fixed by
construction, the curve grades the metric and never a prediction, which is what
makes it the place to ask whether an arbitrary scoring convention changes a
score. Six metrics are run on it below.

\paragraph{One convention, disclosed and then dropped} A cell sentence lists
only expressed genes, so a rank-based metric needs a convention for the unranked
panel genes: \texttt{worst} gives them the last rank $G$, \texttt{middle} a
fixed middle rank $G/2$. Spike-in $\DEdr$ reads $0.952$ against $0.954$ and
nothing here moves by more than $0.01$, so all figures are quoted under
\texttt{worst}. That diagnostic went to no artifact, so it is a check that was
run, not a reproducible result.

\paragraph{The metric this investigation relies on instead} $\NIR$ is Wu
et al.'s ``transposed-rank'' metric from PerturBench~\cite{perturbench}, which
Miller et al.~\cite{miller} identify as one of the few consistently calibrated
metrics across 14 Perturb-seq datasets. It asks not how closely a prediction
matches its own truth, but whether it matches its own truth \emph{more closely
than it matches anyone else's} --- much harder to answer by accident.

\begin{defn}{Normalised Inverse Rank ($\NIR$)}
For a prediction $\hat y$ of condition $i$, with truths $\{y_j\}$ over a declared
comparison set $J$,
\begin{equation}
 \NIR(\hat{y}, i) \;=\; \frac{1}{|J|}\sum_{j \in J}
 \Big[\mathbf{1}\{\sigma(\hat{y}, y_i) > \sigma(\hat{y}, y_j)\}
 + \tfrac{1}{2}\,\mathbf{1}\{\sigma(\hat{y}, y_i) = \sigma(\hat{y}, y_j)\}\Big],
\end{equation}
the fraction of the other conditions in $J$ whose truth the prediction resembles
less than its own, ties counted at one half. $\sigma$ is a \emph{similarity}, so
higher means closer; one built from a distance is that distance negated. The
score counts comparisons, so it reads $\sigma$ only through the ordering that
$\sigma$ induces on $J$, and no strictly increasing transform of $\sigma$ changes
it. Only $\sigma$ and $J$ change across frames, both declared per frame in
\cref{tab:sigma}, with the single exception the table's note records: the
token-retrieval row is computed with a strict inequality and no tie term.
\end{defn}

\Cref{fig:nir} draws one comparison set under this definition, with the three
reading rules that follow from it, the half-tie rule among them.

% fig-nir.tex -- NIR as a discrimination score. Goes with the metric definition
% in Sections/04a-ruler.tex, and is \input from there.
%
% No data dependency: this is a schematic. Its one numeral, NIR = 6/8, is read
% off the drawing itself -- six of the eight hollow circles lie left of the
% filled dot -- and agrees with the definition and the tie-term paragraph in
% 04a-ruler.tex. The sibling fig-nir.json records that check.
%
% This is the v2 COPY of thesis/figs/fig-nir.tex. The v1 file is frozen and was
% not touched. Two changes, and nothing else.
%
% (1) MEASURE. The v1 drawing was 302.79pt wide against v2's 404.03pt (142mm)
%     block -- 75% of the measure, so it sat as a small island with 36mm of
%     white at its right. It is redrawn at the full measure: every x-coordinate
%     is v1's times 1.598, and the marks that ARE the drawing (dot radii, band
%     height, brace amplitude, arrowhead) scale with them, so the figure is
%     enlarged rather than stretched. Type is NOT scaled: \scriptsize and \small
%     stay at the document's own sizes. That is the whole point -- v1's figures
%     were scaled by LaTeX and their type shrank with them. The vertical rhythm
%     of the three reading rules opens from 0.34cm to 0.42cm, proper leading for
%     8pt type rather than the 9.7pt v1 had.
%
%     The explicit bounding box below belongs to this change. A TikZ bounding
%     box includes each node's invisible inner sep, so a picture whose BOX is
%     flush with the measure has its INK sitting 3.5mm inside it at both edges,
%     and the figure reads as not quite reaching the margin. Declaring the box
%     to be the ink fixes that. Measured on the built page: ink 141.82mm inside
%     the 142.00mm block, 0.31pt clear of the left margin and 0.21pt of the
%     right. Nothing is scaled in LaTeX -- no \resizebox, no scale= key, no
%     width= key -- and the document builds with zero overfull boxes.
%
% (2) LEADER. "own truth" floated above the filled dot with a hollow circle
%     1.4cm to its right, and at this width a reader can take it to be labelling
%     that hollow circle. A 0.35cm leader in quietink now ties the label to the
%     dot it names. A plain line, not an arrow: it identifies, it does not
%     assert a direction.
%
% NOT changed: the nine circles and their irregular spacing, NIR = 6/8, all
% three reading rules including the half-tie rule (\S4.7 turns on the half-tie
% property, so its in-figure statement has to survive), and the caption.
\begin{figure}[htbp]
% MARGINALIA (06-figures.md): wider than the 124mm text column, so it
% sits in fullwidth; the caption below spans the text column only.
\begin{fullwidth}\raggedright
\begin{tikzpicture}[
  x=1cm, y=1cm, font=\small,
  lbl/.style={font=\scriptsize, text=ink!70},
  strong/.style={font=\scriptsize, text=ink!90},
  own/.style={circle, fill=ink!85, inner sep=3.2pt},
  other/.style={circle, draw=ink!55, fill=white, inner sep=3.05pt},
  ar/.style={-{Stealth[length=5.5pt]}, ink!65},
  leader/.style={quietink, line width=0.5pt},
]

% The box is the INK, not the ink plus every node's invisible inner sep: left is
% the axis line's left end, right is the right edge of the NIR label's glyphs,
% and the vertical pair reproduces the natural bounding box. 14.190cm in a
% 14.200cm block. Declared FIRST, so nothing drawn below can enlarge it; every
% mark in the picture sits inside it, so nothing spills into the margin either.
\path[use as bounding box] (0.713,-3.080) rectangle (14.903,1.290);

\node[lbl, anchor=south] at (3.28,0.74) {other drugs in the same comparison set};
\node[strong, anchor=south] at (10.07,0.74) {own truth};
% change (2): the leader. 0.35cm, quiet ink, label's south down to the dot's edge.
\draw[leader] (10.07,0.74) -- (10.07,0.39);

\fill[ink!10, rounded corners=1.5pt] (0.88,0.00) rectangle (10.07,0.44);
\foreach \x in {1.52, 2.80, 4.15, 5.67, 7.11, 8.71} {\node[other] at (\x,0.22) {};}
\node[own] at (10.07,0.22) {};
\node[other] at (11.43,0.22) {};
\node[other] at (12.54,0.22) {};
\node[anchor=west, font=\small] at (13.325,0.22) {$\NIR = \tfrac{6}{8}$};

\draw[ar] (0.72,-0.34) -- (13.02,-0.34);
\node[lbl, anchor=north] at (6.87,-0.48)
  {similarity of the prediction to each condition's truth $\longrightarrow$};

\draw[decorate, decoration={brace, mirror, amplitude=4.5pt}, ink!55]
  (0.88,-1.12) -- (10.07,-1.12);
\node[strong, anchor=north] at (5.43,-1.30) {the shaded fraction \emph{is} the score};

\node[lbl, anchor=west] at (0.72,-1.96) {own truth ranked first $\Rightarrow \NIR \to 1$};
\node[lbl, anchor=west] at (0.72,-2.38) {own truth ranked at random $\Rightarrow \NIR = 0.5$};
\node[strong, anchor=west] at (0.72,-2.80)
  {every similarity \emph{identical} $\Rightarrow \NIR = 0.5$ by the half-tie rule, not $0$};

\end{tikzpicture}%
% The comment character above is load-bearing. Without it the line break is an
% interword space inside the centred line, which spends 3.4pt of a measure the
% drawing now fills exactly and pushes the drawing 1.7pt off centre.
\caption[NIR is a discrimination score]{\captiontitle{$\NIR$ asks whether a prediction resembles its own
condition's truth more than it resembles the truths of the other conditions in the comparison set
$J$.} Own truth ranked first gives $\NIR \to 1$; ranked at random gives $0.5$; and a predictor whose
similarities are all identical --- the zero vector, which is what ``predict the average drug
response'' becomes in residual space --- scores $0.5$ by the half-tie rule, not $0$.}
\label{fig:nir}
\end{fullwidth}
\end{figure}


\paragraph{Why chance is $0.50$, and what that claim rests on} Chance sits at
$0.50$ under a symmetry assumption, not as a consequence of using the same truth
set, which is necessary but insufficient. Three things must hold for an
uninformative prediction: $y_i$ exchangeable with the comparison truths
$\{y_j\}$; its label conferring no privileged relation to the prediction; and
symmetric tie handling. The reference is assumed, not proved, so it is checked
empirically in each frame. In the expression frame the drug-agnostic linear map
and control-copy score $0.500$ and $0.504$; in the residual frame: random,
control-copy and \texttt{orth} score $0.501$, $0.506$ and $0.501$. In the
residual frame one asymmetry survives \emph{between} arms: the split-half
reference is scored against half-sample truths, every baseline against
full-sample truths,\seealso{\cref{sec:methods-data} lists every reference in the
chapter.} which makes the comparison conservative. In the expression frame every
arm is scored against the same half-sample truths.

\paragraph{The tie term is not cosmetic} Discard it as bookkeeping and the
metric breaks in residual space (\cref{sec:methods-residual}), where a predictor
of the generic response is exactly the zero vector, so every comparison is a
tie. A strict inequality scores it $0.000$ instead of the
$0.500$ it deserves, below every random predictor.

\paragraph{The comparison set must be declared too} Conditions are grouped by
(cell line, plate) in the expression frame and by cell line in the residual
frame, and the two are not interchangeable. Widening $J$ across plates moves a
drug-agnostic baseline from $0.161$ to $0.399$ (\cref{sec:res-blind}), more than
double the score for the same predictor. And because the generic
response sits in every truth in equal measure, it cancels, leaving only the
drug-specific difference to decide the score --- the component the reversion
artefact exploits.

\paragraph{The expression frame compares two reconstructions} The model emits a
cell sentence, which is gene ranks and nothing else, so no distance in expression
units can be taken on what it produces. The prediction and the truth it is scored
against are both decoded to an expression vector before any distance is
taken,\seealso{\cref{sec:lit-c2s} states the map and the reconstruction argument
C2S makes with it} through the functional form of the map in
\cref{sec:lit-c2s}, fitted once for the whole run as a single slope and
intercept on the base-$10$ logarithm of a gene's position in the sentence,
clamped at zero, with every panel gene the sentence omits set to zero. The
expression frame's $\sigma$ is therefore a distance between two reconstructions
and not between two measured profiles, and the clamp and the zero fill make it
one the underlying rank profiles do not determine.

\paragraph{Three similarity functions, five declared pairs} Three similarity
functions are computed and stored: a negated Euclidean distance, a correlation
between rank profiles, and a cosine. Naming the function does not name the
metric. What fixes the metric is $\sigma$ composed with the representation it
reads, and the negated Euclidean is computed on three different vectors here.
\Cref{tab:sigma} declares the pair $(\sigma, J)$ behind every quoted figure.


\paragraph{Why each frame takes the similarity it does} 
In the expression frame the number attached to a gene is a statement about where
that gene sits in the ordering. The decode reads a gene's position in the
sentence and nothing else, so a large value means near the front and a small one
means far down, and magnitude is not a scale sitting on top of the signal but
part of what the vector says. Two profiles agreeing in shape while differing in
scale are therefore making different claims --- one puts its genes near the top
of the list, the other puts the same genes in the middle --- and a cosine, which
divides the scale out, would call them the same. The truths in a comparison set
really do differ this way, listing different numbers of genes and so differing in
magnitude as well as in pattern, and the Euclidean distance keeps both. What it
keeps is a fact about the representation rather than about abundance, which is
the price of scoring reconstructions.

\paragraph{In the residual frame the choice does not arise} The expression frame
demanded a judgement. Magnitude carries something there, so keeping it or
discarding it is a decision with a cost either way, and the paragraph above
argues for keeping it. In the residual frame there is no such judgement to make,
and not because the cost happens to be small. The two candidate similarities
return the same number.

What makes that possible is that the residual frame never compares raw residuals.
Both arguments pass through $\phi$ first, and $\phi$ assigns its weights by rank
within a block rather than by anything about the condition: the $r$-th most
up-regulated gene receives $1/\log_2(r+1)$ whichever gene it happens to be, the
$r$-th most down-regulated receives its negative, and every other gene receives
zero. Two conditions whose blocks are full therefore hold the same multiset of
values and differ only in which gene holds which. The encoding has standardised
the magnitude before any comparison takes place, so every truth in the comparison
set has the same norm --- and a quantity that is constant across $J$ cannot
separate one member of $J$ from another. There is nothing in the length for a
cosine to throw away.

That alone makes the cosine harmless. The equality is stronger than harmlessness,
and it follows in two lines. Write $\hat{u} = \phi(\hat{y})$ and $u_j =
\phi(y_j)$, and let $\lVert u_j \rVert = C$ for the scored truth and every truth
in $J$, with $\hat{u} \neq 0$. Expanding the squared distance,
\begin{equation}
 -\lVert \hat{u} - u_j \rVert^{2}
 \;=\; \underbrace{-\lVert \hat{u} \rVert^{2} - C^{2}}_{\text{independent of } j}
 \;+\; 2\,\langle \hat{u},\, u_j \rangle ,
\end{equation}
and because $t \mapsto t^{2}$ is strictly increasing on $t \geq 0$, ordering $J$
by $\sigma_{\mathrm{euc}}$ is ordering it by $-\lVert \hat{u} - u_j \rVert^{2}$,
hence by $\langle \hat{u}, u_j \rangle$ alone. The cosine gives
\begin{equation}
 \sigma_{\cos}(\hat{u},\, u_j)
 \;=\; \frac{\langle \hat{u},\, u_j \rangle}
            {\lVert \hat{u} \rVert \, \lVert u_j \rVert}
 \;=\; \frac{1}{C \lVert \hat{u} \rVert}\,\langle \hat{u},\, u_j \rangle ,
\end{equation}
a positive constant multiple of the same inner product, so it orders $J$ the same
way and breaks the same ties. The two similarities do not take equal values ---
one is a distance and the other is bounded in $[-1, 1]$ --- but $\NIR$ reads
$\sigma$ only through the ordering it induces, and on that they are
indistinguishable. Reporting the residual frame under either yields the identical
score.

The equal-norm hypothesis is what carries the result, and it is a property of the
encoding rather than of the biology: it holds wherever both blocks are full, and a
residual too sparse to fill one is the only case in which the two could come
apart. That is also why the expression frame's argument does not transfer.
There the truths genuinely differ in norm, because they list different numbers of
genes, so magnitude is available to separate them and discarding it costs
something real.

\begin{table}[htbp]
% MARGINALIA: full width -- five text columns whose longest cells run to three
% lines even at 182mm; the caption is above at the text width, so it and the
% one note keep a readable measure. No highlight: no one cell carries it.
\caption{\captiontitle{Every $(\sigma, J)$ pair behind a quoted $\NIR$ figure,
declared once.} A negated Euclidean distance appears in three of the five rows
and reads a different vector in each, so what fixes a metric here is the row and
not the similarity function.}
\label{tab:sigma}
\begin{fullwidth}
\begin{threeparttable}
% Name and sigma are sized so "audit computes it" and "negated Euclidean" each
% hold one line; Where quoted does not hyphenate, so no \cref name splits.
\begin{tabularx}{\linewidth}{@{}L{26.5mm} L{28mm} Z{1.05} Z{0.98} >{\hyphenpenalty=10000}Z{0.97}@{}}
\thead{Name} & \thead{$\sigma$} & \thead{Vectors compared}
  & \thead{Comparison set $J$} & \thead{Where quoted} \\
\midrule
\arm{nir_expr} & negated Euclidean & reconstructed expression pseudobulk,
  prediction and truth both decoded & by (cell line, plate), and by cell line in
  the cross-plate arm of \cref{tab:plateleak}
  & the blindness verdict (\cref{sec:res-blind}), \cref{tab:plateleak},
  \cref{fig:difficulty} \\
\addlinespace[10pt]
$\NIR$ as the $\DRF$ audit computes it & negated Euclidean & measured pseudobulk
  expression, no decode & by cell line, and by (cell line, plate) in the
  same-plate arm & \cref{tab:drf} \\
\addlinespace[10pt]
token retrieval & negated Euclidean & undecoded log-expression, as full, shift
  and residual profiles & the group's other drugs, query and truth disjoint
  halves of one condition & the retrieval rows of \cref{tab:ceilings} and
  \cref{tab:divergence} \\
\addlinespace[10pt]
\arm{nir_rank} & correlation of rank profiles & rank profiles over the
  truth's expressed genes & by (cell line, plate) & the admissibility screen
  below, and nothing else \\
\addlinespace[10pt]
residual frame & cosine & signed rank-discounted residual vectors & by cell line
  & the residual $\DRF$ below, and every residual-frame figure
  (\cref{sec:methods-residual}) \\
\end{tabularx}
\begin{tablenotes}[flushleft]
\item The token-retrieval row is the one departure from the definition above:
its code counts strict wins over the comparison set and carries no half-tie
term. Exact ties between Euclidean distances on continuous vectors are
vanishingly rare, so no figure quoted from that row moves under the tie term,
but the row is not an instance of the equation as written.
\end{tablenotes}
\end{threeparttable}
\end{fullwidth}
\end{table}


\paragraph{Grading the metrics, not the models} The framework is Miller
et al.'s~\cite{miller}. Ask not which model scores highest but which
\emph{metric} can see quality at all. No model enters the audit; every predictor
in it is built from measured expression.

\begin{defn}{Dynamic Range Fraction ($\DRF$)}
A metric is graded by where it places a predictor of known quality on its own
scale. Fix three reference predictors and write $m(\cdot)$ for the metric's value
on each:
\begin{equation}
 \DRF(m) \;=\; \frac{m(\textrm{pos}) - m(\textrm{neg})}%
 {m(\textrm{perfect}) - m(\textrm{neg})} .
\end{equation}
The denominator is the entire range the metric has to work with, from a predictor
carrying no drug information to one that is exactly right. The numerator is how
much of that range a real but imperfect signal occupies. A $\DRF$ near $1$ is a
metric that scores genuine signal almost as highly as perfection, $0$ one that
cannot separate it from nothing at all, and a negative value one that placed the
uninformative baseline \emph{above} the real signal.
\end{defn}

\paragraph{The three reference predictors} All three are pseudobulk profiles of
measured expression, stratified by cell line. For each (drug, cell line) the
cells are split in half: $S_{\textrm{GT}}$ supplies the target every predictor is
scored against, and $S_{\textrm{TD}}$ supplies a technical duplicate of the same
condition, disjoint in cells.

\emph{Perfect} is $S_{\textrm{GT}}$ itself, which fixes the top of the scale at
whatever value the metric returns on an identical pair rather than assuming it is
one.

\emph{Negative} is the drug-agnostic mean, the leave-one-out average over the
other drugs in that cell line. It is a real expression profile and not noise,
which is the point: a metric that rewards it is rewarding what every drug does,
not what this drug did.

\emph{Positive} is the \emph{interpolated duplicate}, the technical duplicate
pulled gene by gene toward the negative control in proportion to how weakly that
gene distinguishes the drug. Per gene,
\begin{equation}
 \mu_{\textrm{ID}} \;=\; \alpha\,\mu_{\textrm{TD}}
 \;+\; (1-\alpha)\,\mu_{\textrm{mean}} ,
 \qquad \alpha \;=\; 1 - p_{\textrm{DEG}} ,
\end{equation}
where $p_{\textrm{DEG}}$ is the gene's differential-expression $p$-value for this
drug against the other drugs in the same cell line, computed on the duplicate
half. A gene that separates the drug from its neighbours keeps its measured
value; a gene that does not is pulled back to what any drug would have produced.
What results is a predictor carrying real drug signal with the noise-only genes
damped --- imperfect, but unambiguously better than knowing nothing, which is
exactly the predictor a working metric must rank above the negative control.
Every metric in \cref{tab:drf} sees the same three profiles.

\paragraph{Reading a negative value as inversion} Taking a negative $\DRF$ to
mean the metric is inverted takes one further step, that the sign belongs to the
metric and not to the run.\gloss{inversion}{A metric placing the drug-agnostic
baseline above the positive control, as a property of the metric and not of one
run.} Three of the four negatives survive it; the fourth is carved out where it
appears.

\paragraph{How the number is formed} $\DRF$ is a ratio whose numerator and
denominator are both estimated from the same data, and that makes three things
decisions rather than defaults: how the point estimate is formed, how it is given
an interval, and against what it is tested.\corrmark{The $\DRF$ bootstrap
resampled the wrong unit and imposed no null; both repaired and both arms rerun
before any figure here was read.}

The estimate is a ratio of means and not a mean of ratios. Each of the three
reference predictors is scored on every (drug, cell line) row, the three scores
are averaged across rows, and the ratio is formed once from those averages.
Forming a ratio per row and averaging afterwards would let any row whose
denominator fell near zero dominate the result. Within the rank-based metrics
ties take average ranks, since otherwise a tied gene's rank would be settled by
its position in the panel, which is a fact about how the panel was ordered and
not about the cell.

The interval is a percentile cluster bootstrap resampling whole cell lines, the
unit that repeats across rows, and it recomputes the entire ratio inside every
draw. That second point is the one that matters. The denominator
$m(\textrm{perfect}) - m(\textrm{neg})$ is not a known constant but an estimate,
so an interval that resamples the numerator while holding the denominator at its
observed value describes the uncertainty in the top half of the ratio and calls
it the uncertainty in the ratio. It comes out too narrow. Letting the denominator
move with the numerator keeps the interval about the quantity actually reported.
The cluster is the cell line while the grouping unit is the (cell line, plate)
pair, and the two are reported separately because they are different things; only
the cell-line channel is clustered, so these intervals ignore the crossed well
dimension and are narrower than fully crossed ones, a cost the declared exception
below states in full.

The test imposes its null rather than assuming it. A $p$ read straight off the
observed resampling distribution asks how often the data disagree with
themselves, which for a clearly positive statistic is never, so the $p$ pins to
the Monte Carlo floor whatever the truth is and nothing is tested at all. The
distribution is therefore recentred on $\DRF = 0$ before anything is counted, and
the $p$ is the fraction of null draws reaching the observed value; with
$\num{2000}$ draws the smallest one attainable is $1/2001$. Holm control runs
across the five metrics at once, since they are tested together and read as one
family.
\paragraph{What the number survives} \Cref{tab:drf} carries the values. Every
choice above guards against the estimator flattering itself; none of them says
whether the configuration that produced the estimate is doing the work, and that
is a separate question with separate answers. Under the imposed null no draw of
$\num{2000}$ reaches $\NIR$'s observed value, so it sits at the floor in both
arms and its estimate stands some forty cluster-robust standard errors from zero
--- but a distance from zero is only as good as the setup that generated it.

Three checks bound the setup. Restricting swap partners to the same plate changes
no sign and matches the arms at $27$ cell lines within plate against $25$ across,
so the separation is not an artefact of comparing across batches. A stricter
configuration, capped at $25$ \emph{groups} and admitting only $18$ distinct cell
lines, still gives $\NIR$ \ci{+0.519}{+0.479}{+0.546} over $335$ rows, so it is
not an artefact of breadth either. And the inversion of the four rivals belongs
to the data rather than to a leak: a leave-one-out correction barely moves the
numbers, real cell lines carrying 20--204 drugs, and it survives holding the
plate constant.

What the argument then uses is the sign pattern and nothing finer, drawn for the
same-plate arm in \cref{fig:calibration}, panel b. No ranking among the negative
metrics is claimed.

\begin{table}[htbp]
% MARGINALIA: full width, not in-column: the caption and the two notes would
% run to some 36 lines of the 48mm margin, past the 30-line limit. Caption
% above at the text width, notes below in the sans. The one highlight is
% NIR's same-plate DRF: the arm that holds the plate fixed, and the one the
% argument draws (fig:calibration, panel b).
\caption{\captiontitle{Under the hardened protocol, $\NIR$ is the only one of the five
whose $\DRF$ is positive, and the sign pattern is identical in both arms.} All
four rivals score the drug-agnostic baseline above the positive control and
exclude zero in both arms. The second column is what separates $\NIR$.}
\label{tab:drf}
\begin{fullwidth}
\begin{threeparttable}
\begin{tabularx}{\linewidth}{@{}Z{1.15} Z{1.0} r r Z{0.85}@{}}
\thead{Metric} & \thead{Scored against} & \thead{$\DRF$ all plates}
  & \thead{$\DRF$ same plate} & \thead{Reading} \\
\midrule
$\NIR$ (measured expression) & other conditions' truths & $+0.635$ & \key{+0.545}
  & the only positive \\
\addlinespace[10pt]
\arm{weighted_r2} & its own truth & $-0.163$ & $-0.104$
  & no directional claim \\
$\paneltau$ & its own truth & $-0.319$ & $-0.250$ & inverted \\
\arm{spearman_expr} & its own truth & $-0.650$ & $-0.500$ & inverted \\
$\DEdr$ & its own truth & $-0.686$ & $-0.515$ & inverted; the default \\
\end{tabularx}
\begin{tablenotes}[flushleft]
\item The $\NIR$ row is the audit's own construction, a negated Euclidean
distance on measured pseudobulk expression with no decode, which is the second
row of \cref{tab:sigma} and not the \texttt{nir\_expr} of the blindness verdict.
\item The all-plates run has $\num{1820}$ rows (one drug within one cell line)
clustered on the $25$ cell lines they nest in; the same-plate run has
$\num{586}$ rows over $44$ cell-line$\times$plate groups spanning $27$ cell
lines, clustered on cell line. $\NIR$ reads \ci{+0.635}{+0.604}{+0.663} and
\ci{+0.545}{+0.516}{+0.569}, Holm $p = 0.002$ at the resampling floor in both
arms. The four negatives are point estimates, each excluding zero on the same
clustering in both arms; \texttt{weighted\_r2} sits closest to zero and its sign
tracks cell count, hence no directional claim.
\end{tablenotes}
\end{threeparttable}
\end{fullwidth}
\end{table}


\begin{figure}[htbp]
\begin{fullwidth}
\raggedright
% drawn at 142mm by thesis_v2/figs/calibration.py, the width of the measure; no
% width= key, so it is placed at scale 1.0 and its type prints at the point
% sizes it was set in.
\includegraphics{fig-calibration}
\caption[A zero-information predictor outscores the model]{\captiontitle{A predictor
carrying nothing about drug, cell or biology scores the field's standard metric
above both the within-well reference and the model, and of five metrics only
$\NIR$ grades in the direction its name implies.} (a) $\DEdr$, $K = 50$,
Pearson, tier 2 unseen drugs, \texttt{worst} convention, $n = 400$ conditions;
the axis runs from $0$ and is not truncated. The grey band marks the within-well
split-half precision reference at $0.76$, the two-real-cell positive control of
\cref{tab:de-artifact}: two halves of the same well, a precision figure and not
a biological replicate. Its width is a drawing device, not an interval, and the
arm marks are bare --- the source carries a normal-approximation interval over
conditions, and this document draws only the clustered bootstrap, so nothing is
drawn rather than an interval of another kind. (b) $\DRF$ for the five metrics
of \cref{tab:drf}, the same-plate arm only: $586$ drugs over $44$
cell-line$\times$plate groups spanning $27$ cell lines, the pale block at each
arrowhead the $95\%$ percentile bootstrap interval clustered on cell line
quoted there.}
\label{fig:calibration}
\end{fullwidth}
\end{figure}

\paragraph{What this replicates, and where it does not} The methodological core
of Miller et al.\ replicates, and this audit sharpens it. Metric choice
dominates, and the
control-referenced correlation they flag is the worst performer in both arms
($\DEdr$, $-0.686$ across plates and $-0.515$ within). Their calibrated set does
not transfer intact. \texttt{weighted\_r2}, this project's counterpart of the
weighted $R^2\Delta$~\cite{mejia} they endorse, is negative in both audited
arms, though closest to zero and sign-tracking cell count, so it carries no
directional claim; the inversions of $\paneltau$ and \texttt{spearman\_expr}
belong to neither calibrated set. Their endorsed $\NIR$ does survive. Being
rank-based is not what protects it; refusing to score anything the generic
response can supply is.

\paragraph{Adopting the instrument is also a defence} Miller et al.\ conclude
that under calibrated metrics deep models outperform uninformative baselines;
\cref{sec:res-blind} finds ours at chance on that metric, in the decoded
representation the verdict deploys and not in the one the audit
graded.\seealso{\cref{sec:lit-dispute} sets out the dispute this audit enters.}
Genetic against drug perturbation and a discriminative against a generative
model are candidate explanations; \cref{sec:res-lookup} supplies a third.
Adopting their metric also puts the verdict beyond the objections their own
paper raises.

\paragraph{One declared exception on the intervals} $\DRF$ is a ratio of means;
the two-way estimator used everywhere else here is written for a
mean.\bench{two-way cluster-robust intervals on every
mean}{\cref{sec:methods-data}} Both $\DRF$ intervals are therefore
\emph{one-way} percentile cluster bootstraps that resample whole cell lines and
recompute the ratio inside each draw, respecting the cell-line dependence and
ignoring the crossed well dimension, so they are narrower than fully crossed
intervals. The separation they grade is far too large for that to change the
sign, but the caveat belongs beside the number.

\paragraph{Two questions about the instrument remain, and they differ} The first
is whether the conditions are empirically discriminable at the tested budget. A
spike-in forced choice answers it. Within a plate, a reference pseudobulk from
drug A's cells must be scored closer to a disjoint sample of A's own cells than
to drug B's, chance $0.50$, with contamination of B by A titrating the
difficulty. At zero contamination three of the six metrics run on it
separate the two populations essentially perfectly ($\paneltau = 1.000$,
$\DEdr = 0.995$, top-$n$ $\tau = 0.987$, plate held constant); at full
contamination those same three return $0.501$--$0.502$. The whole-profile
similarities are much weaker even at zero contamination, cosine $0.740$,
Spearman $0.680$, cosine on the shift $0.620$, which matters because a
whole-profile similarity is what $\NIR$ uses here. Split-half retrieval
establishes discriminability, not predictability from untreated input, so
discriminability is not the sole bottleneck at this budget.

The spike-in cannot answer the second question, whether $\NIR$ registers a true
positive, since $\NIR$ was not among the metrics run on it. The evidence is the
$\DRF$ audit's own positive control. Under $\NIR$ on measured expression the
interpolated
duplicate scores $0.802$ against the drug-agnostic mean's $0.457$ across plates,
and $0.668$ against $0.272$ within plate; in the residual frame a split-half
replicate reaches $0.854$ against negatives at $0.501$--$0.506$. None of those is
the reference the verdict is measured against. That reference is the within-well
split-half precision estimate recomputed on the tier-2 within-plate benchmark,
where it sits above chance but not far above it; that, and not $1.0$, is the bar
\cref{sec:res-blind} states the model against.

\begin{scopelimit}[carried into \cref{sec:res-blind} and the rest of the chapter]
The audit licenses $\NIR$ and nothing else. $\paneltau$ is inverted in both
audited arms, so no absolute claim rests on it and no ordering read off it counts
as evidence of skill; an inverted metric preserves neither. Where $\DEdr$
appears later it is only as a paired comparison of two arms on one scale, never
as an absolute score. And a metric is not calibrated by having a relative
calibrated. What is graded here is a negated Euclidean distance on measured
pseudobulk expression, the second row of \cref{tab:sigma}, and the residual
frame's cosine similarity in its own frame, each against its own controls. The
\texttt{nir\_expr} of the blindness verdict is that distance read on decoded
pseudobulk, the first row, with controls of its own but no $\DRF$ of its own, so
it stands on the assumption above. A third $\sigma$ would need its own audit.
\end{scopelimit}

\begin{claim}
This audit grades five candidates, and $\NIR$ is the only
one with a positive $\DRF$ (\ci{+0.635}{+0.604}{+0.663} on a one-way cluster
bootstrap over the $25$ cell lines, the declared exception to this chapter's
two-way rule, Holm $p = 0.002$; \ci{+0.545}{+0.516}{+0.569} over $27$ cell lines
with same-plate partners, Holm $p = 0.002$), and the other four are negative with
intervals excluding zero in both audited arms. Three of them ($\paneltau$,
\texttt{spearman\_expr} and $\DEdr$) are inverted rather than merely
uninformative, each ranking the drug-agnostic baseline above the positive
control; \texttt{weighted\_r2} sits closest to zero and its sign tracks cell
count, so no directional claim is made for it. The saturating predictor is not a
curiosity. It carries no information about drug, cell or biology and scores
$\DEdr = 0.9999$ in rank space, above the two-real-cell positive control of
$0.76$ and above the model's $0.73$. Under $\NIR$ on measured expression the positive control is retrieved well clear of
the negative ($0.802$ against $0.457$ across plates, $0.668$ against $0.272$
within), so the instrument registers a true positive. Every drug-use claim in
the rest of this chapter that rests on a continuous prediction score is
therefore made on $\NIR$, in the audited construction where the audit reaches
and in the decoded one that inherits its licence where it does not; $\DEdr$
appears below only where two arms are compared with each other on the same
scale, never as an absolute score of the model's drug use. The remaining instruments (forced-choice grading,
output invariance, and the logit-space activation tests of
\cref{sec:res-mechanism}) are deliberately on scales of their own, each with a
fixed chance value or a paired control.
\end{claim}

% ===========================================================================
% Beat 2.
% ===========================================================================

\beatsection{E}{2}{In the expression frame, the model is drug-blind}{sec:res-blind}
\label{sec:methods-controls}   % legacy anchor, so v1 cross-references resolve


\begin{question}[2]
With a properly calibrated metric in hand, the investigation can ask the
question it exists to answer. Does the model use the drug at all to make predictions? Four
instruments answer it, and each exists because one specific objection would
otherwise be fatal.
\end{question}

\noindent
The answer is no, and reaching it takes more than a scoreboard. With the metric
calibrated, the model can finally be scored on the tiers, and two readings follow
that look nothing alike and settle nothing on their own.

Within plate on tier 2 ($n = \num{606}$) the model scores $0.498$ against a
within-well split-half precision reference of $0.576$.\bench{one estimator,
seven settings}{\cref{tab:ceilings}} The gap is narrow, but so is the scale: the
reference itself stands only a little above chance, and the model sits at chance
inside it, matched by predictors that know nothing about the drug --- a
drug-agnostic linear map at $0.500$ and a control-copy at $0.504$, against the
global mean's $0.180$. That average also hides a split. Binning the same $\num{606}$
conditions by their own reference gives $296$ on which the reference falls
\emph{below} chance, at $0.287$, where not even a second sample of the same well
scores above chance; $106$ marginal; and $204$ identifiable, whose reference
clears the pre-set $0.80$ bar. On that last third the model looks far better and
discriminates well above chance, reaching $0.768$. A control-copy carrying no
drug information whatsoever reaches $0.766$.

Neither reading is causal, and that is the whole difficulty. A model scoring at
chance might be ignoring the drug, or reading it and failing to turn it into
anything; a model scoring $0.768$ where a drug-blind predictor scores $0.766$
might be adding nothing, or adding something too small for the metric to resolve.
The two readings even point opposite ways, one suggesting the model has learned
nothing and the other that it has learned a good deal, and both are consistent
with the drug token never being consulted. No comparison of scores can separate
these, because no comparison of scores intervenes on anything. What settles it is
an experiment in which the drug changes and nothing else does.

\paragraph{Four instruments, and each closes a different objection} The four are
not four attempts at one measurement. The second is the experiment the previous
paragraph asked for, and the others exist because that experiment is still scored
by something, and whatever scores it could be blamed for the result. Each
successive instrument removes another part of the apparatus from the argument.

The first is the benchmark itself, run \emph{within plate}. Drug and plate are
partially confounded by the atlas's design, so a comparison set free to span
plates lets a batch effect stand in for drug identity and inflates every arm at
once. Restricting the set to a single plate removes that channel before anything
else is asked.

The second is the scramble ablation, and it is the only manipulation among the
four. The control cell is held fixed and only the drug name and mechanism in the
prompt are replaced, by a different drug on the same plate; everything the model
sees apart from the drug is unchanged. A gap between the two generations is
caused by the drug or by nothing, which is what no comparison of scores could
establish. Here the partner is drawn at random, and that carries a known risk: a
swap to a drug with a similar true response makes an unchanged output correct
rather than blind, which biases the test toward finding nothing. A swap-distance
sweep bins (drug, swap) pairs by the response distance between the real and the
swapped-in drug, turning the single test into a curve, and
\cref{sec:res-generalisation} returns to the question with partners stratified by
residual cosine, where a neutral null can be identified rather than assumed.

The third is forced-choice grading. A prediction is scored against its own
condition's truth and against another drug's truth from the same comparison set,
and what is recorded is how often its own wins. 

The fourth is output invariance, and it touches no truth at all. It asks only
whether changing the drug token changes the generated text, comparing generations
against other generations from the same fixed control cell. No similarity
function against a ground truth, no comparison set and no decode enters it, so no
scoring choice can confound it. It is the least informative of the four and the
least confoundable, and those are the same property.

Taken together they share a model and little else. The first can be defeated by
batch leakage, the second by a badly chosen swap partner, the third by a
degenerate comparison set, the fourth by nothing except the model genuinely
ignoring the drug --- and it is that lack of a common failure mode, rather than
any one of them, that carries the verdict.

\paragraph{One measurement rule comes first} Drug and plate are partially
confounded by the atlas's design
(\cref{sec:methods-data}),\framecue{\Eframe}{all four instruments are scored in
the expression frame} so grouping by cell line alone lets the batch effect
stand in for drug identity. A drug-agnostic mean baseline, containing no drug
information whatsoever, scores $0.399$ when the comparison set may span plates
against $0.161$ when it may not, and the precision reference falls from $0.625$
to $0.575$ alongside it, the batch
effect having inflated that too.\seealso{\cref{tab:plateleak} is the leak
audit, a wider run than the benchmark} That is plate leakage, measured and not
assumed. Being wider, its within-plate figures sit slightly off the ones above;
what it licenses is the rule, not a substitute set of scores.

\begin{scopelimit}[carried into every discrimination result in this chapter]
Every discrimination instrument gained a \texttt{--same\_plate\_only} mode and
every discrimination result in this chapter was re-run within plate, except
where a departure from that rule is named at the point of use, as with the
swap-distance sweep that answers the fourth objection below, whose swap partners
were not plate-restricted (\cref{sec:app-sweep}), and the residual-frame scores
of \cref{sec:methods-residual}, whose comparison set is grouped by cell line and
whose justification is given there. A third departure is the identifiability
curve in \cref{tab:ceilings}: it is cross-plate because no within-plate sweep
reaches the cell budgets it reports, the plate-matched version stopping at
twenty cells per condition, and it is computed on raw pseudobulk rather than on
the decoded expression the rest of that table uses. It is therefore read as a
statement about aggregation and never set beside the within-plate references
above it. Conclusions are unchanged; some magnitudes shrink.
\end{scopelimit}

\paragraph{The central manipulation is the scramble ablation} The prompt is left
byte-identical except that the drug name and mechanism are replaced by another
drug's, the control cell and the scoring truth unchanged. Control- and batch-derived leakage therefore cancels in
the scramble gap $\gap = \NIR_{\textrm{model}} - \NIR_{\textrm{scramble}}$, the
only difference in this chapter whose leakage cancels by construction. It does
not displace the drug-agnostic comparisons above and below, which rest on the
within-plate rule instead.

Where that ablation is applied is a choice, and this chapter makes the hard one
first. Every scramble in this section is scored on \texttt{tier2}, held-out
drugs the checkpoint never saw in training, and the substituted partner is drawn
from that same held-out tier --- so the model is asked to distinguish two drugs
it knows nothing about. Nothing here is measured on trained conditions; the
\texttt{tier1} arm of these runs returned no scoreable conditions at all, since
$\NIR$ needs a per-condition pseudobulk and that tier averages roughly two cells
per condition. \Cref{sec:res-encoding} later runs the easier version of the same
manipulation, scoring a model on conditions it did train on, where a drug-blind
model has memorisation to fall back on and any real drug sensitivity has its
best chance of showing. That the expression-frame arms fail here, on the harder
test, is what this section establishes; that the easier test does not rescue
them is \cref{sec:res-epsilon}'s result, and that one target construction does
separate under it is \cref{sec:res-encoding}'s.

% PLACEMENT ONLY. The float was anchored six paragraphs lower (just before "The
% third instrument"), which is where its [htbp] first became satisfiable -- the
% top of the NEXT page -- and that page then had no room left for
% \cref{tab:blind}, which was pushed a further page past the sentence that
% introduces it (measured: anchor p35, caption p37). Anchoring the figure at the
% paragraph it actually depicts, the scramble ablation itself, is earlier in the
% flow, so the figure takes the top slot one page sooner and the table follows on
% the page after its own anchor. Nothing about the figure or its caption changes.
%
% The paragraph above now carries the \cref, so the anchor is explicit and not
% merely positional. Its second sentence lost the restatement "Both arms share
% the same control cell, so ..." because the figure now sets that clause on one
% unbroken line beneath the two arms; the inference is carried by "therefore"
% off "the control cell ... unchanged" in the sentence before it. No claim,
% hedge or magnitude changed, and the paragraph is one line shorter than it was.
%
% \input, not \includegraphics: figs/fig-scramble.tex is TikZ and carries its
% own float, caption and \label. \input resolves against the main-document
% directory, so this is thesis_v2/figs/, NOT the read-only ../thesis/figs/ that
% \graphicspath adds for raster/PDF assets only.
\begin{figure}[htbp]
% MARGINALIA (06-figures.md): fits the text column; caption in the margin.
\begin{lrbox}{\tvfb}%
  \begin{tikzpicture}[x=21mm, y=1mm, font=\small]
    \def\ur{180}                     % upper strip, mm per unit of the gap
    \def\lz{-52}\def\ls{63}          % lower strip, origin and mm per unit
    \def\lchance{\lz+(0.5-0.326)*\ls}% $\NIR = 0.50$ on the lower strip
    % Four of the five gaps sit within a few thousandths of zero. The dot stays at
    % its true height and the label drops straight down to one aligned row, so the
    % series reads as the flat line it is rather than as five scattered labels.
    \newcommand{\gapdot}[2]{%
      \fill[ink] (#1,{#2*\ur}) circle (1.5pt);
      \draw[rulegrey, line width=0.3pt] (#1,{#2*\ur}) -- (#1,-12.4);
      \node[anchor=north, font=\scriptsize, text=ink] at (#1,-12.4) {$#2$};}
    \newcommand{\nirdot}[2]{%
      \fill[quietink] (#1,{\lz+(#2-0.326)*\ls}) circle (1.3pt);
      \node[anchor=west, font=\scriptsize, text=quietink, fill=white, inner sep=1pt]
        at (#1+0.06,{\lz+(#2-0.326)*\ls}) {$#2$};}
  
    % === upper strip, the scramble gap ========================================
    \node[anchor=west, font=\scriptsize\bfseries, text=accent]
      at (-0.42,29) {$\gap$, expression frame};
    \node[anchor=east, font=\scriptsize\itshape, text=quietink]
      at (4.40,29) {trend $\mathrm{corr}(\textrm{distance},\gap) = +0.051$};
    \draw[ink, line width=0.5pt, densely dotted] (-0.42,0) -- (4.40,0);
    \node[anchor=west, font=\scriptsize\itshape, text=ink] at (4.42,0) {zero};
  
    % the one bin for which an interval is quoted, drawn as an interval
    \draw[accent!55, line width=1.6pt] (1,{-0.002*\ur}) -- (1,{0.124*\ur});
    \fill[accent] (1,{0.060*\ur}) circle (1.6pt);
    \node[anchor=west, font=\scriptsize, text=accent]
      at (1.08,{0.060*\ur}) {$+0.060$};
  
    \gapdot{0}{-0.046}
    \gapdot{2}{-0.006}
    \gapdot{3}{-0.004}
    \gapdot{4}{-0.019}
  
    % === the divider IS the argument: two quantities, two scales, one x =======
    \draw[rulegrey, line width=0.4pt] (-0.42,-20) -- (4.40,-20);
  
    % === lower strip, the model's own score, which drifts =====================
    \node[anchor=north west, font=\scriptsize\bfseries, text=quietink]
      at (-0.42,-23) {the model's own $\NIR$ in the same bins, on its own scale};
    \draw[ink, line width=0.5pt, densely dotted]
      (-0.42,{\lchance}) -- (4.40,{\lchance});
    \node[anchor=west, font=\scriptsize\itshape, text=ink]
      at (4.42,{\lchance}) {chance};
    \nirdot{0}{0.497}
    \nirdot{1}{0.612}
    \nirdot{2}{0.523}
    \nirdot{3}{0.401}
    \nirdot{4}{0.326}
  
    % === the one axis the two strips share ====================================
    \foreach \i/\lab/\rng in {0/nearest/{2.75--4.03}, 1/{}/{4.03--4.60},
                              2/{}/{4.60--4.93}, 3/{}/{4.93--5.32},
                              4/farthest/{5.32--6.25}}{
      \node[anchor=north, font=\scriptsize, text=quietink, align=center]
        at (\i,-58) {\i\\\lab};
      \node[anchor=north, font=\tiny, text=quietink] at (\i,-66) {\rng};}
    \node[anchor=north, font=\scriptsize\itshape, text=quietink]
      at (2,-72) {swap distance $\lVert y_A - y_B \rVert$, real drug to swapped-in drug};
  \end{tikzpicture}
  %
\end{lrbox}
\sidecap[The scramble gap does not grow with swap distance]{figure}{\captiontitle{Naming a
  drug whose response is far from the truth damages the prediction no more than
  naming a near neighbour.} Results taken in the Tier 2. Upper strip, the scramble gap across five bins of swap
  distance. A real dissimilarity effect would climb from left to right and peak at
  the farthest bin; instead the correlation between distance and gap is $+0.051$,
  and the only bin whose interval is quoted, \ci{+0.060}{-0.002}{+0.124}, is the
  second of five. Lower strip, the model's own $\NIR$ in those same bins, which
  does fall overall, though not monotonically, because drugs far from everything
  populate the far bins. The two strips are different quantities on different
  scales, and $\gap$ uses the same drug on both arms, so that drift cancels out of
  it. Non-canonical: \texttt{max\_new\_tokens=1200}, $n = 40$ per bin, which
  excludes a large dissimilarity effect without establishing a null.\label{fig:sweepbins}}
  
  
\end{figure}

\paragraph{Two constraints on the swap partner} The partner must be a different
drug. Allowing the same drug at another dose or plate substitutes the drug for
itself, an unchanged prediction is then correct and not blind, and the contamination
concentrates in the most similar stratum. It must also be on the same plate,
since residuals retain plate structure and a cross-plate partner would differ in
batch as well as in drug.

A scramble that swaps in a drug with a similar true response is a weak test, and
that risk is structural here. In full-profile space the mechanism annotation is
close to uninformative about transcriptional response. Same-mechanism drugs in
Tahoe are just over two per cent closer to each other than arbitrary pairs
are.\seealso{\cref{sec:res-encoding} quotes the distances behind that statement}
Hence conditioning on mechanism fails to move $\DEdr$ in the grouping ladder of
\cref{sec:res-lookup}, and a partner of different mechanism cannot be relied on
for a different response. Mechanism does carry signal in the residual frame
(\cref{sec:res-scope}); in full-profile space the generic response swamps it, a
frame difference and not a contradiction.

Two remedies are used, one here and one later. The first turns the single test
into a curve. Within a (cell line, plate) group, each drug's cells are split in
half and the distance between two drugs is taken between their half-sample truth
pseudobulks, $\lVert y_A - y_B \rVert$. Because both profiles sit against the
same plate-matched control, that distance is the same whether it is measured on
profiles or on shifts, and it is exactly the distance $\NIR$ already uses to rank
one condition's truth against another's. The sweep therefore inherits the
benchmark's own geometry rather than importing a notion of similarity from
elsewhere, and the axis is fixed before any generation: it is a property of the
atlas, not of the model.

Partners are drawn to span the quantiles of that distance from $A$, so each real
drug contributes swaps from across its whole available range instead of one
partner of unknown similarity. Every swap rebuilds $A$'s prompt with $B$'s name
and mechanism, leaves the control cell and the rest of the prompt intact, and
scores the generation against \emph{$A$'s} truth. The target never moves; only
the label in the prompt does. Pooling the pairs and binning them by distance
gives five bins running from nearest to farthest
(\cref{fig:sweepbins}, \cref{sec:app-sweep}), and the gap is aggregated inside
each bin with intervals clustered on cell line.

What makes the curve worth reading is that each outcome was given its meaning
before the run. A gap rising with distance would say the model does discriminate
dissimilar drugs and that the flat aggregate had been diluted by near-twin swaps
--- a correction to the claim, not a confirmation of it. A gap that stays flat
into the far tail says the blindness holds at the hardest setting the plate
affords, where the swapped-in drug's real response is as unlike the truth as the
data allows.

\paragraph{The third instrument is forced-choice grading} A prediction is scored
against its own condition's truth and against another drug's truth from the same
comparison set, recording how often its own wins. Chance is $0.50$ by
construction, and the precision reference lands between $0.67$ and $0.83$
depending on cells per condition.

\paragraph{The fourth is output invariance} It asks whether the drug token
changes the generated text at all, without reference to any truth. It is run at
$3800$ tokens, the longest budget in the chapter, because
truncation would cut the comparison
itself.\seealso{\cref{sec:methods-data} records the four generation budgets} For
a fixed control cell it compares
% \resizebox, not a reflow: the two \underbrace labels are the width, and both
% are wanted on one line so the subtraction reads as one sentence. Natural
% width is 165.6mm against the 142mm block (67.15pt over), so this sets at
% about 0.86. Nothing here is body text.
\begin{equation}
\resizebox{\linewidth}{!}{$\displaystyle
 \textrm{invariance gap} \;=\;
 \underbrace{\overline{\textrm{sim}}(\textrm{same prompt, resampled})}%
   _{\textrm{sampling noise alone}}
 \;-\; \underbrace{\overline{\textrm{sim}}(\textrm{prompts differing in the
   drug})}_{\textrm{noise plus any drug effect}},
$}
\end{equation}
with similarity as top-$100$ Kendall $\tau$ and as Jaccard overlap, over $8$
contexts $\times$ $12$ drugs at temperature $0.8$. It is the least informative
of the four and the least confoundable.

\begin{table}[htbp]
% MARGINALIA: full width; caption above at the text width, the one note below.
% No highlight: four instruments agreeing is not one cell.
\caption{\captiontitle{Four independent instruments agree, and the two that return
intervals bound the drug's effect instead of merely covering zero.} Near-ceiling
$\DEdr$ is achieved identically with a wrong-mechanism drug, so that score is
drug-independent.}
\label{tab:blind}
% Column 3 was `l', then a weighted X: at an earlier, narrower measure its
% natural width plus columns 1--2 left the last column too narrow for the
% single unbreakable word "indistinguishable", which ran 43.91pt past the
% fullwidth block, and the X let the \ci intervals break after the estimate.
% Column 3 is `l' again so no \ci ever breaks mid-cell; at the 182mm measure
% the first three columns' natural widths leave the Reading column, a Y, about
% 44mm, so "indistinguishable" fits.
\begin{fullwidth}
\begin{threeparttable}
\begin{tabularx}{\linewidth}{@{}l l l Y@{}}
\thead{Instrument} & \thead{Quantity} & \thead{Value} & \thead{Reading} \\
\midrule
Forced-choice grading & own-truth win rate
  & $0.48$ ($\approx$ scramble; reference 0.67--0.83) & chance \\
\addlinespace[10pt]
Output invariance & invariance gap, top-$100$ $\tau$ & \ci{0.000}{-0.019}{+0.016}
  & bounded, $\leq 0.016$ \\
 & invariance gap, Jaccard & \ci{+0.000}{-0.008}{+0.008}
  & bounded, $\leq 0.008$ \\
\addlinespace[10pt]
Scramble, tier 1 & $\DEdr$ & $0.740$ real vs $0.739$ scrambled
  & indistinguishable \\
Scramble, tier 2 & $\DEdr$ & $0.733$ real vs $0.729$ scrambled
  & indistinguishable \\
Scramble (identifiable subset) & $\gap$ ($\NIR$) & \ci{+0.014}{-0.016}{+0.042}
  & bounded, $\leq +0.042$ \\
\end{tabularx}
\begin{tablenotes}[flushleft]
\item The two $\DEdr$ rows are paired comparisons on one scale, the only use
\cref{sec:res-metrics} licenses for that metric.
\end{tablenotes}
\end{threeparttable}
\end{fullwidth}
\end{table}

% PLACEMENT ONLY. This is one of the three survivors the \paragraph orphan patch
% in preamble.tex section 4 records: the global reservation is
% 3\baselineskip, which this head cleared with 2.3 lines to spare and then set
% itself plus exactly two lines at the foot of the page. 5 lines is head plus
% four, which the page cannot fake.
\needspace{5\baselineskip}%
\paragraph{An interval that covers zero earns its keep by what it excludes} Not
by the word \emph{null}. Two predictions from a byte-identical prompt agree at
$\tau = 0.198$, so the invariance interval caps any drug-driven divergence at
roughly a twelfth of that agreement, and the Jaccard version is tighter still,
\ci{+0.000}{-0.008}{+0.008}. On the identifiable stratum the scramble gap of
\ci{+0.014}{-0.016}{+0.042} runs against the $0.170$ separating the model's
$0.768$ from that stratum's $0.938$ reference; at the very top of the interval,
naming the right drug would buy a quarter of what is there to be bought, and the
point estimate puts it at $8\%$ of that. Neither figure shows the effect is
zero, and no experiment here could. Both put it far below the headroom these
conditions leave.


\begin{figure}[htbp]
\begin{fullwidth}
  \raggedright
  % drawn at 142mm by thesis_v2/figs/difficulty.py; no width= key, so it is placed
  % at scale 1.0 and its type prints at the point sizes it was set in.
  \includegraphics{fig-difficulty}
  \caption[Difficulty structure and the stratum ladder]{\captiontitle{Conditions are
  more empirically discriminable at the tested pseudobulk budget where the
  precision reference is high, and precisely there, a predictor containing no drug
  information matches the model.} Both panels are $\NIR$ in the expression frame,
  on one shared axis. (a) The within-well split-half precision reference against
  cells per condition, for two comparison sets: cross-plate, and within plate. (b)
  The same reference used to bin conditions, five arms per stratum, $n = 296$
  unwinnable, $106$ marginal, $204$ identifiable. On the unwinnable stratum the
  reference itself is $0.287$, below chance, as are the model, control-copy and
  scramble marks. On the identifiable stratum the model reaches $0.768$ and a
  control-copy baseline $0.766$, paired contrast \ci{+0.002}{-0.026}{+0.028},
  clustered over cell lines. Plotting the model alone would suggest it does better
  on higher-reference drugs. The arm marks are bare: the source carries no
  interval for a stratum mean, and this document draws only the clustered
  bootstrap, so nothing is drawn rather than an interval of another kind.}
  \label{fig:difficulty}
\end{fullwidth}
\end{figure}
  \paragraph{Four ways the verdict could be an artefact} Each is a way the result
  could be about the apparatus rather than about the model, and each is settled by
  a measurement rather than by argument.
  
  \emph{That there is nothing to learn.} A ladder of mean-shift baselines, each
  predicting the average shift of a progressively narrower group, says
  otherwise.\seealso{\cref{sec:res-lookup} gives the ladder} On tier 2 the global
  mean scores $0.056$ and conditioning on cell line raises it about
  two-and-a-half-fold to $0.142$, so the atlas carries real structure. What it does
  not carry is structure the drug label reaches: mechanism leaves the global figure
  untouched at $0.052$, and mechanism together with cell line does no better than
  cell line alone ($0.120$ against $0.142$). There is something to know, and the
  variable that knows it is the cellular context.
  
  \emph{That the benchmark is too weak to say anything.} Barely half the atlas
  clears the pre-set identifiability bar (\cref{tab:scale}), so the verdict is
  re-scored where the wells are demonstrably adequate. It does not survive the
  move. On that stratum the model appears to beat the linear baseline, $0.768$
  against $0.531$, and then a control-copy carrying no drug information reaches
  $0.766$ on the same conditions; the paired contrast is
  \ci{+0.002}{-0.026}{+0.028} and the leak-immune $\gap$ stays bounded as above.
  
  That control-copy carries a second finding, and it is easy to walk past. Across
  the three strata its score tracks the reference almost step for step
  (\cref{fig:difficulty}, panel b): $0.310$ against a reference of $0.287$ where
  the reference is lowest, $0.540$ against $0.689$ on the marginal stratum, $0.766$
  against $0.938$ on the identifiable one. A predictor that knows nothing about any
  drug does best exactly where the ranking problem is easiest. So what makes a
  condition identifiable is largely what makes its untreated cells distinctive
  among their plate-mates, and the distance from the model to a high reference is
  not all room that better drug modelling would fill. The reference does rise with
  aggregation (\cref{fig:difficulty}, panel a: $0.61$ at $10$ cells, $0.76$ at
  $40$, $0.85$ at $120$); what the strata add is that most of the ranking can be
  done without the drug at all.
  
  \emph{That the encoding discards the drug.} A cell sentence keeps the rank order
  of genes and throws away expression magnitude, so if the drug lived in the
  magnitudes the following sections would be attacking the wrong object. This needs
  no model to test: ask whether true expression discriminates drugs better than
  rank does. It does not. At single-cell resolution, cosine, Pearson and Spearman
  similarities on true normalised expression all sit within $0.04$ of chance
  (\ci{0.529\textrm{--}0.531}{0.515}{0.543}), the same near-chance regime as rank.
  What recovers the signal is aggregation: at fifteen cells the control-referenced
  shift cosine reaches \ci{0.793}{0.775}{0.809}. The re-encoding of
  \cref{sec:res-encoding} will therefore not be a restoration of magnitudes; it
  keeps rank throughout and changes only what is ranked. The same measurement also
  bounds what a single cell can deliver, since the single-cell same-drug against
  different-drug gap is only \ci{+0.006}{+0.003}{+0.008}: these instruments do not
  resolve drug identity from one observed cell at this sequencing depth and
  $946$-gene panel.
  
  \emph{That the scramble swaps in near-twins.} The swap-distance sweep bounds this
  rather than settling it. It sits outside the within-plate rule, an earlier run at
  \texttt{max\_new\_tokens=1200} on the unseen-drug tier of a different evaluation
  set, with swap partners not plate-restricted and $n = 40$ per bin --- wide enough
  to exclude a large effect, not to establish a null (\cref{sec:app-sweep}). No
  distance bin has an interval excluding zero. The far tail reads $-0.019$ across a
  $2.3\times$ spread in swap distance with a trend correlation of $+0.051$, so a
  large dissimilarity effect is excluded for the expression-frame model and a small
  one is left open.
  
  \paragraph{What none of the four can see} All four watch the output. They
  establish that the generated text does not move when the drug does, and that is
  the strongest behavioural statement available. It is also silent on why. A model
  that never encodes the drug and a model that encodes it and uses it incorrectly
  produce the same flat instruments, and no experiment that only reads generations
  can separate them. The ladder above already hints at where to look, since the
  asymmetry it finds between cellular context and drug class is one the activations
  turn out to repeat. \Cref{sec:methods-probe} therefore stops asking what the
  model emits and starts asking what it holds.

    \paragraph{A fifth reading of the same sweep, with no comparator in it} 
    The four
  instruments above each contrast the model against something --- a baseline, a
  control-copy, a scrambled prompt --- and each is only as good as the thing it is
  set against. The swap-distance sweep supports one further reading that contrasts
  nothing. Its partners were drawn to span the distance quantiles, five per
  condition, so within a single condition the same model answers the same question
  five times with a differently wrong drug named each time. Write $y_{ik}$ for the
  $\NIR$ of condition $i$'s $k$th scrambled prompt and $x_{ik}$ for the
  \emph{negated} Euclidean distance between that condition's truth and the truth of
  the drug substituted into it --- negated so that $x$ rises with similarity, as
  the partner cosine does everywhere else in this thesis, and one predicted sign
  serves every version of the test. Then
  \begin{equation}
   y_{ik} \;=\; \alpha_i \;+\; \beta\,x_{ik} \;+\; \varepsilon_{ik},
   \qquad k = 1,\dots,5,
  \end{equation}
  with one intercept $\alpha_i$ per condition. The target-drug prompt is never a
  point on this line, so nothing is measured against a chosen baseline, and
  $\beta = 0$ is the null by construction: a model that ignores the name emits the
  same distribution whichever drug is substituted, putting its five points at one
  height whatever their similarities happen to be. Prompt sensitivity predicts
  $\beta > 0$, the score rising as the substituted drug's real response moves
  closer to the target's.

  Which intercepts are fitted decides the answer. A single pooled intercept draws
  one line through all $200$ points at once, so conditions are compared with each
  other as much as with themselves, and fitted that way the slope is $+0.0986$
  $[+0.0180, +0.1792]$ clustered on drug --- a clear effect, and not a real one. A
  condition whose drug is unlike anything else on its plate has distant partners
  \emph{and} an unusual truth that any prompt will struggle to rank, so the line
  bends without a name having been read. Giving each condition its own intercept
  removes that comparison and leaves only what varies inside a condition: one
  model, one target, five differently wrong drugs. Almost the whole effect is in
  the part that discards --- the between-condition slope is $+0.1563$ against a
  within-condition slope of $+0.0055$, $[-0.0472, +0.0582]$ clustered on drug and
  $[-0.0484, +0.0594]$ on cell line, over $200$ prompts across $40$ conditions.
  The tighter interval excludes the pooled estimate it replaces. Compared with
  itself, no condition's score moves with how unlike the substituted drug is.

  What this adds is not more of the same evidence. The scramble is a paired
  intervention --- one condition, one thing changed --- and is confound-free by
  construction, but it needs a comparator, and \cref{sec:res-generalisation} shows
  how much of a gap the choice of comparator can carry. This regression needs none.
  What it gives up in exchange is what it can license: every point in it is a
  \emph{wrong} prompt, so it can show only whether the model tells wrong answers
  apart, never whether it produces a right one. A model could pass it and predict
  nothing. Here that trade costs nothing, because the model does not pass it.

\begin{claim}
The model is at chance overall ($0.498$ against a $0.576$ within-well split-half
precision reference), matched on the identifiable stratum by a
zero-drug-information control-copy ($0.768$ against $0.766$, paired contrast
\ci{+0.002}{-0.026}{+0.028}), and no more moved by a maximally dissimilar swap
than by a similar one in the single-cell column of a non-canonical
\texttt{max\_new\_tokens=1200} sweep outside the within-plate rule
(\cref{sec:app-sweep}). 
\end{claim}
